\documentclass[11pt]{article}

\usepackage[margin=0.7in]{geometry}
\usepackage{amsmath,amssymb}
\usepackage{booktabs}
\usepackage{array}
\usepackage{tabularx}
\usepackage{enumitem}
\usepackage[colorlinks=true,linkcolor=blue,citecolor=blue,urlcolor=blue]{hyperref}
\usepackage{titlesec}
\usepackage{parskip}

\titleformat{\section}{\normalfont\large\bfseries}{\thesection}{0.6em}{}
\titlespacing*{\section}{0pt}{1.2ex}{0.4ex}
\setlist{nosep,leftmargin=1.2em}
\renewcommand{\arraystretch}{1.05}
\setlength{\parskip}{0.5ex}

\newcommand{\R}{\mathbb{R}}

\begin{document}

\begin{center}
{\Large\bfseries Configuration Snapshot}
\end{center}

\vspace{-1ex}

% =====================================================================
\section{Pipeline}

\begin{small}
\begin{tabularx}{\linewidth}{@{}lX@{}}
\toprule
Hardware    & Unitree Go2 + Livox Mid-360 LiDAR; ROS 2 Humble on Jetson Orin. Lab compute: RTX 5090. \\
Sim         & Isaac Lab on \texttt{Isaac-Velocity-Flat-Unitree-Go2-v0}. \\
Locomotion  & Off-the-shelf flat-trained policy, frozen black box. Drop-in safety filter --- works with any user's locomotion stack. \\
Robot model & Single-integrator $\dot{x}_{\text{robot}} = u$; $L_g h \neq 0$ so standard CBF applies. \\
Outer loop  & RL teacher emits robust-CBF parameters $(\alpha,\varphi,a,c)$ at 50 Hz; closed-form half-space QP outputs $u_{\text{safe}} \in \R^{3}$ ($v_x, v_y, \omega_{\text{yaw}}$); locomotion maps that to 12-D joint targets; output to \texttt{walking\_bridge}. \\
\bottomrule
\end{tabularx}
\end{small}

% =====================================================================
\section{CBF formulation and four-parameter action space}

Safe-set $h(x)$ is built from per-shape analytical SDFs (cylinders, post-v2.12
commitment) combined via Eq.~19 multi-obstacle SDF + Eq.~20 exponential
smoothing $h(p) = \lambda(1 - \exp(-\gamma \cdot \text{sdf}(p)))$; robot footprint
enters as Minkowski expansion by $0.15$ m. The QP is solved closed-form on GPU
as a half-space projection. We include the $L_f h$ obstacle-drift term so
$\dot h$ is correctly accounted for when obstacles move. With
$\partial h / \partial p_{\text{obs}} = -L_g h$ by symmetry,
$\dot h = L_g h \cdot u - L_g h \cdot v_{\text{obs}} = L_g h \cdot (u -
v_{\text{obs}})$, so the constraint becomes:
\[
  L_g h \cdot u_{\text{safe}}
  \;\geq\;
  -\alpha\,(h - c) \;+\; \varphi \,\lVert L_g h\rVert^{2} \;+\; a
  \;+\; L_g h \cdot v_{\text{obs}}
\]
The teacher emits a 5-D action $(\alpha,\varphi,a,b,c)$, of which 4 are active
(see table). $b$ is reserved for an input-dependent slack $b\,\lVert u\rVert$
that would require an SOCP solver; closed-form half-space projection breaks
otherwise --- deferred.

\begin{small}
\begin{tabularx}{\linewidth}{@{}llX@{}}
\toprule
Param     & Range          & Role \\
\midrule
$\alpha$  & $[0.1, 5.0]$   & Class-$\mathcal{K}$ slope on shifted $h(x)-c$; absorbs model / tracking error. \\
$\varphi$ & $[0.0, 5.0]$   & Coefficient on $\lVert L_g h\rVert^{2}$ (RHS); absorbs actuation uncertainty. \\
$a$       & $[0.0, 3.0]$   & Additive RHS slack; absorbs state-independent measurement uncertainty. \\
$c$       & $[0.0, 1.0]$   & Inward shift of safe-set boundary; boundary correction when $h$ underestimates. \\
\bottomrule
\end{tabularx}
\end{small}

% =====================================================================
\section{Privileged observation}

The teacher conditions on privileged information available in sim but not at
deploy. A student in Wk3 will reproduce its latent from LiDAR + history.

\begin{small}
\begin{tabularx}{\linewidth}{@{}lX@{}}
\toprule
Total dimension   & \textbf{8207-D} \\
Dynamics block    & 15-D --- proprioception, base velocity, COM offset, mass, friction estimate. \\
Occupancy grid    & 8192-D --- $2 \times 64 \times 64$ ego-centric grid at $0.1$ m resolution, 6.4 m FOV (current frame + previous frame, so the encoder can infer obstacle motion). \\
\bottomrule
\end{tabularx}
\end{small}

% =====================================================================
\section{Network architecture}

\begin{small}
\begin{tabularx}{\linewidth}{@{}lX@{}}
\toprule
Encoder         & priv $\to$ CNN [Conv $2{\to}16$ s$=2$, $16{\to}32$ s$=2$, Linear $\to 64$], dyn MLP [$15 \to 64$] $\to$ concat $\to$ Linear $128 \to 12$. \\
Latent          & $Z \in \R^{12}$ (RMA-style bottleneck; \texttt{get\_z(obs)} exposed for student replay). \\
Policy          & $\pi_{\text{teacher}}: Z \to (\alpha,\varphi,a,b,c)$, $128$ hidden, output dim $5$, tanh-squash to physical-range scale. \\
$u_{\text{des}}$ & Never enters the network --- sidecar to the CBF; only used in reward via $\lVert u_{\text{safe}}-u_{\text{des}}\rVert^{2}$. CBF params should depend on environment, not on user's current command. \\
\bottomrule
\end{tabularx}
\end{small}

% =====================================================================
\section{PPO training recipe}

\begin{small}
\begin{tabularx}{\linewidth}{@{}lX@{}}
\toprule
Algorithm       & PPO via \texttt{rsl\_rl}, on-policy. \\
Optimizer       & AdamW, weight\_decay $= 1\mathrm{e}{-5}$. \\
Entropy coef.\  & $0.005$ (load-bearing --- $0.001$ caused action-std collapse in earlier runs). \\
Action-rate     & $-0.005 \cdot \lVert \Delta a \rVert^{2}$ in the reward (smooths CBF params in time). \\
Smoothness story& Orthogonal mechanisms: weight-decay smooths input$\to$output map; action-rate smooths in time. \\
Episode length  & 20 s. \\
Parallel envs   & $4096$. \\
Iterations      & $5000$ ($\sim$5 h on RTX 5090; SPS $\sim$27 K). \\
Health logging  & 14 CBF stats per iter: 8 = $(\alpha,\varphi,a,c)$ mean/std (slot adaptation); 2 = $h_{\min}, h_{\text{mean}}$ (closeness to obstacle); 2 = QP-active rate, $u_{\text{safe}}$-clamp rate; 2 = noise-$\sigma$ mean/std (DR active). \\
\bottomrule
\end{tabularx}
\end{small}

% =====================================================================
\section{Reward stack}

\begin{small}
\begin{tabularx}{\linewidth}{@{}llX@{}}
\toprule
Term                            & Weight                                                       & Trigger \\
\midrule
\texttt{collision}              & $-100$                                                       & Terminal --- per-shape SDF $< 0$ (obstacle hit). \\
\texttt{base\_contact\_penalty} & $-100$                                                       & Terminal --- robot fall (matches \texttt{collision} symmetrically). \\
\texttt{stuck}                  & $-2.0$ / step                                                & Per-step --- when $\lVert v_{xy}\rVert < 0.15$ m/s. \\
\texttt{infeasibility}          & $-10$ / step                                                 & Per-step --- QP infeasibility marker. \\
\texttt{u\_safe\_deviation}     & $-0.1\cdot\lVert u_{\text{safe}}-u_{\text{des}}\rVert^{2}$   & Per-step --- penalize CBF deflection of nominal command. \\
\texttt{proximity}              & $-0.5\cdot\exp(-\text{min\_sdf}/0.5)$                        & Per-step --- distance-shaped pressure away from obstacles. \\
\texttt{action\_rate}           & $-0.005\cdot\lVert\Delta a\rVert^{2}$                        & Per-step --- smooths $(\alpha,\varphi,a,b,c)$ over time. \\
\bottomrule
\end{tabularx}
\end{small}

\noindent
\textbf{Design principle: DR-implicit parameter shaping.} No reward terms
explicitly target $(\alpha, \varphi, a, c)$ behavior. The DR axes provide the
disturbance; the parameters learn to absorb it. This is the cleaner paper
claim --- ``DR creates the disturbance, the parameter learns to compensate'' ---
rather than hand-crafted reward signals per slot.

% =====================================================================
\section{Domain randomization}

\begin{small}
\begin{tabularx}{\linewidth}{@{}lXll@{}}
\toprule
Axis             & Knob                                  & Train range                      & OOD eval range \\
\midrule
Friction         & static / dynamic                      & $(0.30, 1.20)$ / $(0.20, 1.00)$  & $(0.15, 1.50)$ / $(0.10, 1.30)$ \\
Disturbance      & external force / torque               & $\pm 10$ N / $\pm 2$ Nm          & $\pm 18$ N / $\pm 3.5$ Nm \\
COM offset       & xy / z                                & $\pm 5$ cm / $\pm 3$ cm          & $\pm 8$ cm / $\pm 5$ cm \\
Obstacle motion  & per-episode $\lvert v_{\text{obs}}\rvert$ cap & $\sim\mathcal{U}(0, 0.4)$ m/s & up to $0.6$ m/s \\
Obstacle density & \texttt{separation\_buffer}           & $0.4$ m                          & $0.2$ m \\
Obstacle pool    & \multicolumn{3}{l}{Cylinder-only (SHIELD-style commitment); $K_{\text{actual}} \in [0, 20]$ unique per reset; radii $0.10$--$0.50$ m.} \\
Moving obstacles & \multicolumn{3}{l}{$\sim 50\%$ drift constant velocity per episode (random direction), capped by per-episode $v_{\text{obs}}$.} \\
\bottomrule
\end{tabularx}
\end{small}

\noindent
\textbf{Per-episode persistent perception bias} (load-bearing for the $a$ and
$c$ slots --- the change that distinguishes the current configuration from
earlier versions). At each episode reset, every (env, obstacle) draws:
\[
  \sigma_e \sim \mathcal{U}(0,\,\sigma_{\max}),
  \quad \varepsilon_{e,k} \sim \mathcal{N}(0,\,\sigma_e^{2}\,I_2)
  \;\;(\text{persistent for the episode})
\]
The QP sees biased obstacle positions $p_k^{\text{QP}} = p_k +
\varepsilon_{e,k}$; the priv-obs grid stays clean. Per-episode persistence
(NOT per-step i.i.d.) is critical: i.i.d.\ noise would average to zero under
LLN over 1000-step episodes and leave $a$/$c$ with no gradient signal. The
biased structure matches the temporally-coherent error of a real
LiDAR-cluster perception pipeline: cf.\ SHIELD (Yang et al.\ 2025,
arXiv:2505.11494, Section V), where Livox Mid-360 returns are passed
through Euclidean clustering and each cluster is treated as a cylinder
of fixed radius $R = 0.3$ m, with cluster-centroid drift between LiDAR
frames as the dominant perception-error source. Training
$\sigma_{\max} = 0.05$; OOD env \texttt{Isaac-CBF-Go2-NoisyPerception-v0}
uses $\sigma_{\max} = 0.10$ to test the $a$ slot at deploy stress.

% =====================================================================
\section{Multi-planner training mix}

The teacher trains against a mixture of nominal commands (this is what the
locomotion controller would see at deployment):

\begin{small}
\begin{tabularx}{\linewidth}{@{}lXX@{}}
\toprule
Planner                     & Share & Behavior \\
\midrule
\texttt{smooth\_goal}        & $0.40$ & Smoothed goal-direction velocity command (most common deploy case). \\
\texttt{waypoint}            & $0.25$ & A*-style waypoint chasing through clutter. \\
\texttt{mpc}                 & $0.20$ & MPC-like receding-horizon planner. \\
\texttt{legacy\_goal}        & $0.05$ & Plain goal-direction. \\
\texttt{walk}                & $0.05$ & Constant forward walk (open-space behavior). \\
\texttt{adversarial}         & $0.05$ & Random-direction commands (regularizer). \\
\bottomrule
\end{tabularx}
\end{small}

\noindent
\textbf{Bimodal resample DR.} Each episode rolls $P=0.5 \to$ planner
re-samples every $\mathcal{U}(5,15)$ s (mid-switch); $P=0.5 \to$ planner
locked for the full $100$ s episode. \textbf{Eval is always locked
(deploy-realistic).} The training/eval mismatch is the design point:
training mid-switch teaches intrinsic recovery from CBF deflection; locked
eval verifies that recovery transfers to a single deployed nav stack.

\noindent
\textbf{All planners rate-limited, MPC re-derived (2026-05-08).} Every
planner emits commanded $v_{xy}$ subject to $\lVert \Delta v \rVert \leq
\text{max\_accel} \cdot dt$ ($\text{max\_accel} = 2.0$ m/s$^{2}$), so
nominal commands cannot step-discontinuously and create
artificial-tracking-error stress on the locomotion controller. MPC was
additionally re-tuned ($K_p = 2.5$, $K_d = 0$) after solo-planner
bisection (Section~9) showed its prior PD-for-double-integrator law was
the dominant K{=}0 stuck contributor.

% =====================================================================
\section{Failure-rate layering}

Where do the absolute fall and stuck numbers come from? We isolate by running
a \texttt{--no\_obstacles} diagnostic ($K_{\text{actual}}=0$ at eval), then
bisect further with \texttt{--solo\_planner} to attribute per-planner
contribution. This bisection caught a hidden MPC-planner bug that had been
silently inflating absolute numbers across all v2.6--v2.11 runs.

\begin{small}
\begin{tabularx}{\linewidth}{@{}Xcc@{}}
\toprule
Cumulative layer                                                                   & Fall        & Stuck       \\
\midrule
Native locomotion in its own flat-velocity training task                           & $0.5\%$     & ---         \\
$+$ our env (50 Hz, full DR), $K{=}0$, vanilla \texttt{UniformVelocityCommand}     & $2.9\%$     & $5.7\%$     \\
$+$ multi-planner mix (rate-limited, MPC fixed)                                    & $1.8\%$     & $6.5\%$     \\
$+$ obstacles $+$ plain hand-tuned CBF baseline (B0)                               & $\sim 20\%$ & $\sim 17\%$ \\
\bottomrule
\end{tabularx}
\end{small}

\begin{itemize}
  \item \textbf{$0.5\% \to 2.9\%$ fall jump.} Cost of our DR $+$ 50 Hz command streaming on locomotion alone, before any obstacle exists. This is the locomotion floor.
  \item \textbf{Multi-planner mix adds essentially zero floor.} Combined $8.3\%$ (mix) vs.\ $8.6\%$ (vanilla) is statistically indistinguishable. The 6-planner mix is no longer a confound for absolute numbers.
  \item \textbf{Hidden MPC bug fixed (2026-05-08).} The previously-reported $16.5\%$ stuck floor was \emph{not} a goal-reached artifact (as I had assumed). Solo-planner bisection traced it to an MPC math error: a PD-for-double-integrator control law $u = K_p\,(p_{\text{goal}} - p) - K_d\,v$ with $K_p = 0.5$, $K_d = 0.5$, used as a velocity command (not an acceleration). $K_p = 0.5$ together with the $1.2$ m/s velocity cap produced a $2.4$ m deceleration zone (cap $\div K_p$); the robot crawled for most of every goal-approach. $K_d \neq 0$ added oscillation near goal. Fix: $K_p = 2.5$ (decel zone $\to 0.48$ m), $K_d = 0$, rate-limit on all 6 planners (\texttt{max\_accel} $= 2.0$ m/s$^{2}$). Solo-MPC stuck collapsed $44.4\% \to 5.4\%$; full-mix stuck collapsed $16.9\% \to 6.5\%$.
  \item \textbf{CBF-attributable piece.} The residual once locomotion $+$ DR floors are subtracted --- what the teacher actually targets while preserving the forward-invariance guarantee on contact.
\end{itemize}

% =====================================================================
\section{Evaluation protocol}

\textbf{Baselines.} \;
\textbf{B0:} fixed $(\alpha,\varphi,a,b,c) = (0.5, 0, 0, 0, 0)$ --- plain CBF,
no robust slacks. \;
\textbf{B1:} B0 with tuned $\varphi \in \{0.5, 1.5, 3.0\}$ --- ISSf form. \;
\textbf{B2:} TISSf form, varying $(\alpha, \varepsilon_0, \lambda)$. \;
All baselines hand-tuned over a $24$-config grid; reported = best per task.

\noindent
\textbf{B-fixed-X ablations} (paper Table 2): at eval time only, clamp ONE
slot to a fixed value (mean of training-time output for that slot); the other
3 stay adaptive. Tests whether per-slot adaptation is load-bearing on each OOD
axis: \emph{Bf-$\alpha$} on DensePack/in-dist, \emph{Bf-$\varphi$} on
HighDisturbance, \emph{Bf-a} on NoisyPerception, \emph{Bf-c} on HeavyCOM /
FastObstacles.

\noindent
\textbf{Metric.} Combined fall + stuck rate (lower is better). $64$ envs per
config $\times$ $2000$ sim steps $\approx 150$ episodes per row; same
\texttt{--num\_envs}, \texttt{--steps\_per\_config}, seeded RNG across configs.

% =====================================================================
\section{Headline result --- validated baseline}

Trained teacher vs.\ a 24-config sweep of plain, ISSf, and TISSf hand-tuned
baselines. Margin = BR (best response, ours) $-$ best hand-tuned baseline,
combined fall + stuck (lower is better).

\begin{small}
\begin{tabularx}{\linewidth}{@{}lXl@{}}
\toprule
Eval                & Type                                         & Margin (BR vs.\ best baseline) \\
\midrule
In-distribution     & mixed                                        & \textbf{+6.9 pp WIN} \\
Slippery            & priv-obs (continuous friction shift)         & \textbf{+5.6 pp WIN} \\
HighDisturbance     & priv-obs (episodic force / torque)           & \textbf{+9.0 pp WIN} \\
FastObstacles       & priv-obs (grid history)                      & \textbf{+10.6 pp WIN} \\
HeavyCOM            & priv-obs (startup COM bias)                  & \textbf{+5.9 pp WIN} \\
DensePack           & scene-only ($h(x)$ symmetric)                & $+0.6$ pp tie \\
RealisticCompound   & all 5 modest pushes simultaneously           & $-0.3$ pp tie \\
\bottomrule
\end{tabularx}
\end{small}

\noindent
Average single-axis margin $+6.4$ pp ($5$ wins, $1$ scene-only tie, $1$
compositional tie). Pattern: the teacher wins on every OOD axis where it has
asymmetric privileged information (friction, force, COM, velocity history);
ties where the playing field is symmetric (DensePack --- same $h(x)$ for both
methods) or compositional.

\noindent
\textbf{Status note (updated 2026-05-09).} v2.12 finished training and
the 8-eval landed: $3$W$/1$T$/4$L, average margin $+0.17$ pp. In-dist
($+7.05$), NoisyPerception ($+5.51$), FastObstacles ($+7.05$) are
clean wins; DensePack/Slippery/HighDist/Compound regressed vs v2.6.
Headline finding from B-fixed-X ablations: \textbf{Bf-$a$ on
NoisyPerception = $+44.1$ pp degradation} --- fixing $a$ to its mean
breaks performance on perception-noisy environments by $44$ pp combined.
The $a$ slot is genuinely load-bearing despite a small training-time
mean ($\approx 0.05$); the policy uses it in spike-mode when episode
$\sigma_e$ is large. Strong validation of the four-parameter architecture.

% =====================================================================
\section{Two findings driving v2.14}

The session also surfaced two architecturally-meaningful findings.

\noindent
\textbf{Finding 1 --- per-step $\varphi$ output is jitter, not adaptation.}
B-fixed-$\varphi$ on HighDisturbance INVERTED: clamping $\varphi$ to its
mean BEAT the per-step adaptive policy by $7.6$ pp combined. Diagnosis
(framed by my mathematical advisor's catch): per-step Gaussian sampling
on the $\varphi$ output is the same LLN/CLT-style failure mode we fixed
on the input side for perception-bias DR. PPO's advantage averages
reward over the $\sim 800$-step episode, so any single $\varphi_t$
contributes $\sim 1/800$ to the return $\Rightarrow$ no usable
gradient signal for state-conditional $\varphi$ $\Rightarrow$ policy
converges to "noisy mean." B-fixed-$\varphi$ replicates the mean
without the per-step noise, hence the inversion.

The structural root: $\varphi$ scales the actuation-uncertainty robust
term $\varphi \lVert L_g h \rVert^2$. Actuation uncertainty is an
\textbf{environment-class} property (motor lag, friction, locomotion
tracking error --- stable across an episode), not a state-conditional
one. Per-step adaptation has no useful target.

Symmetric fix: lock $\varphi$ once at episode start, replay for the
remainder of the episode. The Gaussian exploration noise that was on
$\varphi_t$ at reset stays in (now between-episode); the per-step
jitter that was poisoning the constraint is removed. Predicted v2.14
outcome: B-fixed-$\varphi$ on HighDist swaps from $-7.6$ to positive.

\noindent
\textbf{Finding 2 --- the SHIELD-perception gap.} I evaluated the v2.12
checkpoint under the deploy-realistic SHIELD perception pipeline (QP
sees uniform $R = 0.3$ m cylinders for every obstacle, regardless of
true radius --- exactly Yang et al.\ 2025 Sec~V). The BR margin
flipped from $+7.05$ (priv perception) to $-3.5$ pp (SHIELD), a
$10.5$ pp swing. Mechanism: BR's outputs are calibrated for the QP
knowing the true radii; under SHIELD, the QP under-pads big obstacles
($R = 0.50$ m treated as $R = 0.3$ m) producing $+6$ pp falls.
Importantly, the baselines actually IMPROVED under SHIELD ($-7.2$ pp
combined) because their fixed parameters were already calibrated for
"the average obstacle" --- which is what uniform $R = 0.3$ m
approximates.

Implication: the v2.12 architecture's win depends on the QP knowing
true geometry, which is not available at deploy time on real LiDAR.
Need to train UNDER SHIELD perception so the policy learns to
compensate for the $R = 0.3$ m QP assumption.

\noindent
\textbf{v2.14 spec.} Both findings folded into a single training run
launching next ($\sim 5$ h train + $\sim 5$ h eval). Two fixes
relative to v2.12:
\begin{itemize}
  \item \textbf{Per-episode $\varphi$ lock} (capture once, replay; matches the v2.12 perception-bias persistence on the input side).
  \item \textbf{Train under SHIELD-style perception} (\texttt{perception\_mode = "shield\_v0a"}: ground-truth positions, uniform $R = 0.3$ m cylinders).
\end{itemize}

The two fixes are architecturally orthogonal and both have predicted
gain. Successful v2.14 produces a deployment-realistic adaptive CBF
with all four slots load-bearing under realistic perception ---
a stronger result than v2.6 / v2.12 because the win number now
corresponds to the deployment perception model, not the privileged
training one.

\end{document}
